\section{Лекция 5}
Напомним определение запутанных и факторизованных состояний системыю Пусть имеется N гильбертоы=вых пространств $\mathcal{H}_1, ..., \mathcal{H}_N$, пусть $\ket{\psi} \in \mathcal{H} = \mathcal{H}_1\otimes ...\otimes\mathcal{H}_N$. Пусть $\alpha = \{j^{(1)}, ..., j^{(N)}\}$ --- общий индекс, а в $\mathcal{H}$ введен естественный базис $\ket{\eta_\alpha} = \ket{\phi_{j^{(1)}}}\otimes\ket{\kappa_{j^{(2)}}}\otimes...\otimes\ket{\chi_{j^{(N)}}}$. Тогда вектор $\ket{\psi}$ можно разложить по этому базису как $\ket{\psi} = \sum\limits_\alpha C_\alpha \ket{\eta_\alpha}$.\\

\begin{definition}
	Если для любого такого базиса $\{\ket{\eta_\alpha}\}$ получаем, что всегда существует 2 или более ненулевых коэффициента $C_\alpha$, то состояние $\ket{\psi}$ называется \textbf{запутанным}.
\end{definition}
\begin{definition}
	Если существует хотя бы один базис $\{\ket{\eta_\alpha}\}$ такой, что состояние $\ket{\psi}$ представимо в виде $\ket{\psi} = \ket{\phi_{\tilde{j}^{(1)}}}\otimes\ket{\kappa_{\tilde{j}^{(2)}}}\otimes...\otimes\ket{\chi_{\tilde{j}^{(N)}}}$ (т.е. все коэффициенты, кроме одного равны нулю), то состтояние называется \textbf{факторизованным}.
\end{definition}

\begin{example}
	Рассмотрим два спина 1/2: $s^{(1)} = 1/2\in \mathcal{H}_2^{(1)}$, $s^{(2)} = 1/2\in \mathcal{H}_2^{(2)}$. Пусть $\mathcal{H}_4 = \mathcal{H}_2^{(1)}\otimes\mathcal{H}_2^{(2)}$.\\
	
	Напомним, что $\ket{\pm^{(i)}} = \ket{s^{(i)}=1/2,~ s_z^{(i)} = \pm 1/2}$, $i =1,~2$, $\ket{\eta_{++}} = \ket{+^{(1)}}\otimes\ket{+^{(2)}}$ и т.п.\\
	
	Пример факторизованного состояния: $\ket{s=1,~ s_z = +1,~...} = \ket{+^{(1)}}\otimes\ket{+^{(2)}}$. Пример запутанного состояния: $\ket{s=0,~ s_z = 0,~...} = \dfrac{1}{\sqrt{2}}(\ket{+^{(1)}}\otimes\ket{-^{(2)}}-\ket{-^{(1)}}\otimes\ket{+^{(2)}})$, покажем, что это состояние действительно является запутанным.
	\begin{proof}
		От противного. Пусть существует базис, в котором это состояние факторизовано, т.е. представимо в виде $\ket{s=0,~s_z=0,~...} = \ket{\phi^{(1)}}\otimes\ket{\chi^{(2)}}$. Тогда новый базис можно выразить через старый:
		\[
		\begin{cases}
			\ket{\phi^{(1)}} = a^{(1)}\ket{+^{(1)}} + b^{(1)}\ket{-^{(1)}}\\
			\ket{\chi^{(2)}} = a^{(2)}\ket{+^{(2)}} + b^{(2)}\ket{-^{(2)}}\\
		\end{cases}
		\]
		Подставим:
		\[
		\ket{\phi^{(1)}}\otimes\ket{\chi^{(2)}} = a^{(1)}a^{(2)}\ket{+^{(1)}}\otimes\ket{+^{(2)}} + a^{(1)}b^{(2)}\ket{+^{(1)}}\otimes\ket{-^{(2)}} + b^{(1)}a^{(2)}\ket{-^{(1)}}\otimes\ket{+^{(2)}} + b^{(1)}b^{(2)}\ket{-^{(1)}}\otimes\ket{-^{(2)}}
		\]
		Приравняем коэффициенты с исходным состоянием:
		\[
		\begin{cases}
			\begin{cases}
				a^{(1)}a^{(2)} = 0\\
				b^{(1)}b^{(2)} = 0\\
			\end{cases} \Longrightarrow a^{(1)}a^{(2)}b^{(1)}b^{(2)} = 0\\
			\begin{cases}
				a^{(1)}b^{(2)} = \dfrac{1}{\sqrt{2}}\\
				b^{(1)}a^{(2)} = -\dfrac{1}{\sqrt{2}}\\
			\end{cases} \Longrightarrow a^{(1)}a^{(2)}b^{(1)}b^{(2)} = -\dfrac{1}{2}\\
		\end{cases} \Longrightarrow \text{ противоречие.}
		\]
	\end{proof}
\end{example}

Существуют и более сложные примеры запутанных состояний, например состояние \textbf{Гринберга-Хорна-Цайлингера}: $\ket{\psi^{(GHZ)}} \in \mathcal{H}^{(1)}\otimes\mathcal{H}^{(2)}\otimes\mathcal{H}^{(3)}$, $\ket{\psi^{(GHZ)}} = \dfrac{1}{\sqrt{2}}(\ket{+^{(1)}}\otimes\ket{+^{(2)}}\otimes\ket{+^{(3)}} + \ket{-^{(1)}}\otimes\ket{-^{(2)}}\otimes\ket{-^{(3)}})$.\\

Очевидно, что матрица плотности факторизованного состояния $\ket{\psi} =\ket{\phi_{j^{(1)}}}\otimes\ket{\kappa_{j^{(2)}}}\otimes...\otimes\ket{\chi_{j^{(N)}}}$ имеет вид
\[
\hat{\rho} = \ket{\psi}\bra{\psi} = \hat{\rho}_{j^{(1)}}^{(1)} \otimes\hat{\rho}_{j^{(2)}}^{(2)} \otimes ... \otimes\hat{\rho}_{j^{(N)}}^{(N)},
\]
где $\hat{\rho}_{j^{(i)}}^{(i)} = \ket{\phi_{j^{(i)}}^{(i)}}\bra{\phi_{j^{(i)}}^{(i)}} \ \forall i\in\{1, .., N\}$, $\ket{\phi_l^{(i)}}\in\mathcal{H}^{(i)}$.\\

\begin{definition}
	Если существует базис, в котором матрицу плотности состояния $\hat{\rho}$ можно представить в виде суммы матриц плотности факторизованных состояний
	\[
	\hat{\rho} = \sum\limits_l w_l \hat{\rho}_l^{(1)} \otimes ... \otimes\hat{\rho}_l^{(N)},
	\]
	где $\sum\limits_l w_l = 1$, $\hat{\rho}_l^{(i)} = \ket{\phi_l^{(i)}}\bra{\phi_l^{(i)}}$, $\ket{\phi_l^{(i)}}\in\mathcal{H}^{(i)}$, то такое состояние называется \textbf{сепарабельным}.
\end{definition}

Не существует общего критерия, устанавливающего, что данная конкретная матрица плотности описывает именно сепарабельное или именно запутанное состояние.  Такой критерий существует только для двух частных случаев: системы кубит-кубит и системы кубит-кутрит.\\

\begin{theorem}[Необходимое условие сепарабельности для двух квантовых систем произвольной размерности]
	Пусть дана квантовая система, состоящая из двух подсистем. Если ее матрица плотности является сепарабельной, то после частичного транспонирования эта матрица все еще будет удовлетворять определению матрицы плотности. Для матрицы плотности запутанного состояния поведение неопределено.\\
	Если система имеет вид кубит-кубит или кубит-кутрит, то это \textbf{критерий Переса-Городецких}.
	
\end{theorem}
\begin{proof}
	Пусть задана сепарабельная матрица плотности $\hat{\rho} = \sum\limits_l w_l \hat{\rho}_l^{(1)}\otimes\hat{\rho}_l^{(2)}$.\\
	Заметим, что если матрица $\hat{\rho}$ удовлетворяет условиям на матрицу плотности, то и $\hat{\rho}^T$ удовлетворяет тем же условиям (это можно доказать непосредственной проверкой).\\
	Под частичным транспонированием подразумевается транспонирование отдельных матриц плотности в сумме, т.е.
	\[
	\hat{\rho}^{T_2} = \sum\limits_l w_l \hat{\rho}_l^{(1)}\otimes(\hat{\rho}_l^{(2)})^T \ \text{ или } \ \hat{\rho}^{T_1} = \sum\limits_l w_l (\hat{\rho}_l^{(1)})^T\otimes\hat{\rho}_l^{(2)}
	\]
	Очевидно, что это все еще матрицы плотности.
\end{proof}





\subsection{\S X.3 Матрица плотности без классического прибора}
Ранее мы вводили матрицу плотности $\hat{\rho}$ следующим образом: имеется некий классический прибор, измеряющий наблюдаемую величину и с некоторыми вероятностями $w_1,~w_2, ...,~w_N$ выдающий соответственно состояния $\ket{\psi_1},~\ket{\psi_2}, ...,~\ket{\psi_N}$, тогда матрицей плотности назовем $\hat{\rho} = \sum\limits_{l=1}^N w_l\hat{\rho}_l$, где $\hat{\rho}_l = \ket{\psi_l}\bra{\psi_l}$ --- проекторы на чистые состояния. Но на самом деле матрицы плотности возникают в квантовой механике и без участия классического измерительного прибора.\\

Вернемся к примеру $s^{(1)} = 1/2,~ s^{(2)} = 1/2$. Рассмотрим синглет $\ket{s=0,~s_z=0,~...} = \dfrac{1}{\sqrt{2}}(\ket{\eta_{+-}} - \ket{\eta_{-+}})$ --- это чистое состояние. Вопрос: в каком состоянии находится каждая из подсистем? Покажем, что в смешанном.\\
Введем $\hat{\rho} = \ket{s=0,~s_z=0,~...}\bra{s=0,~s_z=0,~...}$. Тогда матрица плотности подсистемы
\[
\hat{\rho}^{(1)} = Tr_{\mathcal{H}^{(2)}}\hat{\rho} = \langle +^{(2)} | \hat{\rho} | +^{(2)} \rangle + \langle -^{(2)} | \hat{\rho} | -^{(2)} \rangle
\]
\[
\begin{cases}
	\braket{+^{(2)}}{s=0,~s_z=0,~...} = \dfrac{1}{\sqrt{2}}(\ket{+^{(1)}}\braket{+^{(2)}}{-^{(2)}} - \ket{-^{(1)}}\braket{+^{(2)}}{+^{(2)}}) = -\dfrac{1}{\sqrt{2}}\ket{-^{(1)}}\\\\
	\braket{-^{(2)}}{s=0,~s_z=0,~...} =\dfrac{1}{\sqrt{2}}\ket{+^{(1)}}\\
\end{cases} \Longrightarrow
\]
\[
\hat{\rho}^{(1)} = \left(-\dfrac{1}{\sqrt{2}}\right)^2\ket{-^{(1)}}\bra{-^{(1)}} + \left(\dfrac{1}{\sqrt{2}}\right)^2\ket{+^{(1)}}\bra{+^{(1)}} = \dfrac{1}{2}(\ket{-^{(1)}}\bra{-^{(1)}} + \ket{+^{(1)}}\bra{+^{(1)}}) = \dfrac{1}{2}\hat{1}
\]
Таким образом, $\hat{\rho}^{(1)}$ --- единичная матрица (с точностью до константы), а единичную матрицу нельзя факторизовать (в смысле квантовых состояний) $\Longrightarrow$ это смешанное состояние. Аналогично, можно показать, что $\hat{\rho}^{(2)} = \dfrac{1}{2}\hat{1}$ --- также смешанное состояние.\\




\subsection{\S X.4 Состояния Белла}
Снова рассмотрим $s^{(1)} = 1/2\in \mathcal{H}_2^{(1)}$, $s^{(2)} = 1/2\in \mathcal{H}_2^{(2)}$, $\mathcal{H}_4 = \mathcal{H}_2^{(1)}\otimes\mathcal{H}_2^{(2)}$.\\

Какие интересные базисы мы можем ввести?

\begin{enumerate}
	\item Естественный базис:
	\[
	\begin{cases}
		\ket{\eta_{++}} = \ket{+^{(1)}}\otimes\ket{+^{(2)}}\\
		\ket{\eta_{+-}} = \ket{+^{(1)}}\otimes\ket{-^{(2)}}\\
		\ket{\eta_{-+}} = \ket{-^{(1)}}\otimes\ket{+^{(2)}}\\
		\ket{\eta_{--}} = \ket{-^{(1)}}\otimes\ket{-^{(2)}}\\
	\end{cases}
	\]
	\item Физический базис: 
	\[
	\begin{cases}
		\ket{s=1,~s_z=+1} = \ket{\eta_{++}}\\
		\ket{s=1,~s_z=0} = \dfrac{1}{\sqrt{2}}(\ket{\eta_{+-}} + \ket{\eta_{-+}})\\
		\ket{s=1,~s_z=-1} = \ket{\eta_{--}}\\
		\ket{s=0,~s_z=0} = \dfrac{1}{\sqrt{2}}(\ket{\eta_{+-}} - \ket{\eta_{-+}})\\
	\end{cases}
	\]
	\item \textbf{Базис Белла}:
	\[
	\begin{cases}
		\ket{\psi_+} = \dfrac{1}{\sqrt{2}}(\ket{+^{(1)}}\otimes\ket{-^{(2)}} + \ket{-^{(1)}\otimes\ket{+^{(2)}}})\\
		\ket{\psi_-} = \dfrac{1}{\sqrt{2}}(\ket{+^{(1)}}\otimes\ket{-^{(2)}} - \ket{-^{(1)}\otimes\ket{+^{(2)}}})\\
		\ket{\phi_+} = \dfrac{1}{\sqrt{2}}(\ket{+^{(1)}}\otimes\ket{+^{(2)}} + \ket{-^{(1)}\otimes\ket{-^{(2)}}})\\
		\ket{\phi_-} = \dfrac{1}{\sqrt{2}}(\ket{+^{(1)}}\otimes\ket{+^{(2)}} - \ket{-^{(1)}\otimes\ket{-^{(2)}}})\\
	\end{cases}
	\]
	Все состояния, состовляющие базис Белла являются запутанными состояниями двух спинов 1/2.
\end{enumerate}




\subsection{\S X.5 Квантовая телепортация}
Имеется произвольное спиновое запутанное состояние $\ket{\psi^{(s)}} = c_1^{(s)}\ket{+^{(s)}} + c_2^{(s)}\ket{-^{(s)}} \in \mathcal{H}_2^{(s)}$, $c_1^{(s)},~c_2^{(s)} \neq 0$. Этот же вектор имеется у Алисы $(A)$, и она хочет передать его без повреждений Бобу $(B)$.\\

Заранее $A$ и $B$ готовят запутанное белловское состояние $\ket{\Phi_{AB}^+} = \dfrac{1}{\sqrt{2}}(\ket{+^{(A)}}\otimes\ket{+^{(B)}} + \ket{-^{(A)}}\otimes\ket{-^{(B)}}) \in \mathcal{H}_{AB}=\mathcal{H}_2^{(A)}\otimes\mathcal{H}_2^{(B)}$, т.е. 1 кубит будет храниться в лаборатории $A$, а другой в лаборатории $B$. Рассмотрим общее гильбертово пространство $\mathcal{H} = \mathcal{H}_2^{(s)}\otimes\mathcal{H}_{AB} = \mathcal{H}_2^{(s)}\otimes\mathcal{H}_2^{(A)}\otimes\mathcal{H}_2^{(B)}$ $\Longrightarrow$ общий вектор состояния
\[
\ket{\psi_{tot}} = \ket{\psi^{(s)}}\otimes\ket{\Phi_{AB}^+} = \dfrac{1}{\sqrt{2}}(c_1^{(s)}\ket{+^{(s)}}\otimes\ket{+^{(A)}}\otimes\ket{+^{(B)}} + c_1^{(s)}\ket{+^{(s)}}\otimes\ket{-^{(A)}}\otimes\ket{-^{(B)}} + 
\]
\[
+ c_2^{(s)}\ket{-^{(s)}}\otimes\ket{+^{(A)}}\otimes\ket{+^{(B)}} + c_2^{(s)}\ket{-^{(s)}}\otimes\ket{-^{(A)}}\otimes\ket{-^{(B)}})
\]
Прямое произведение ассоциативно $\Longrightarrow$ $\mathcal{H} = (\mathcal{H}_2^{(S)}\otimes\mathcal{H}_2^{(A)})\otimes\mathcal{H}_2^{(B)} = \mathcal{H}^{(SA)}\otimes\mathcal{H}^{(B)}$. Здесь $\mathcal{H}^{(SA)}$ относится к тому, что имеется в лаборатории $A$, а $\mathcal{H}^{(B)}$ --- в лаборатории $B$. Разложим по базису Белла в $\mathcal{H}^{(SA)}$:
\[
\ket{\psi_{tot}} = \dfrac{1}{2}[ \ket{\Phi_{AB}^+}\otimes(c_1^{(s)}\ket{+^{(B)}} + c_2^{(s)}\ket{-^{(B)}}) +
\]
\[
+ \ket{\Phi_{AB}^-}\otimes(c_1^{(s)}\ket{+^{(B)}} - c_2^{(s)}\ket{-^{(B)}}) +
\]
\[
+ \ket{\Psi_{AB}^-}\otimes(c_1^{(s)}\ket{-^{(B)}} + c_2^{(s)}\ket{+^{(B)}}) +
\]
\[
\ket{\Psi_{AB}^+}\otimes(c_1^{(s)}\ket{-^{(B)}} - c_2^{(s)}\ket{+^{(B)}})]
\]
Теперь $A$ может измерить состояние в базисе Белла и получить $\ket{\Phi_{AB}^+}$, $\ket{\Phi_{AB}^-}$, $\ket{\Psi_{AB}^+}$ или $\ket{\Psi_{AB}^-}$ (каждое с вероятностью 1/4) и сообщить об этом $B$. Тогда $B$, исходя, из имеющегося у него кубита будет точно знать, какое состояние передала $A$, при этом состояние в лаборатории $A$ уничтожилось.